\chapter{Configurazione completa}
\label{cap:B-configurazione}

\section{Configurazione di piattaforma}

Da \code{platform/control-plane/src/main/resources/application.yml}. Sono omessi
l'elenco \code{spring.autoconfigure.exclude} (le auto-configurazioni di Spring
Boot di cui il control plane non usa alcun bean: client HTTP, TLS, esecutori di
task, multipart, sessioni, AOP, \emph{info} e spazio su disco) e il blocco
\code{management.endpoint}, che abilita solo \code{health} con le probe e
\code{prometheus}. Le chiavi sono raggruppate per argomento; i commenti sono
del paper, non del file.

\begin{lstlisting}[language=yamlnf]
server:
  port: 8080
management:
  server:
    port: 8081
  endpoints:
    web:
      exposure:
        include: health,prometheus

nanofaas:
  metrics:
    profile: basic              # basic | advanced | soak
  admin:
    runtime-config:
      enabled: false
  scheduler:
    strategy: ${NANOFAAS_SCHEDULER_STRATEGY:}   # vuoto = nessuna scelta esplicita
  deployment:
    default-backend: ""         # k8s | container-local; vuoto = risoluzione implicita
    wakeup:
      timeout: 30s
      poll-interval: 250ms
      ready-observation-max-age: 5s
  registry:
    path: "${NANOFAAS_REGISTRY_PATH:build/nanofaas/functions.json}"
  defaults:
    timeoutMs: 30000
    concurrency: 4
    queueSize: 100
    maxRetries: 3
  retry:
    initial-backoff: 100ms
    max-backoff: 2s
  rate:
    maxPerSecond: 1000000
  k8s:
    namespace: ""
    callbackUrl: "http://control-plane.default.svc.cluster.local:8080/v1/internal/executions"
  scaling:
    poll-interval-ms: 5000
    default-min-replicas: 1
    default-max-replicas: 10
  execution-store:
    ttl: 5m
    sync-ttl: 30s
    max-lifetime: 30m
    max-outcomes: 100000
    max-keys: 100000
    max-outcome-bytes: 0        # 0 = max-outcomes x 116 byte
  http-client:
    max-connections: 500
    pending-acquire-max-count: 0        # 0 = 2 x max-connections
    pending-acquire-timeout-ms: 45000
    max-idle-time-ms: 30000
    max-life-time-ms: 0                 # 0 = rotazione disabilitata
    eviction-interval-ms: 5000
    inactive-pool-dispose-interval-ms: 5000
    pool-inactivity-ms: 30000
  container-local:
    proxy:
      max-request-bytes: 2097152
      max-response-bytes: 4194304
      max-buffered-bytes: 33554432
      inbound-read-timeout: 5s
      response-write-timeout: 5s

sync-queue:
  enabled: true
  admission-enabled: true
  max-depth: 200
  max-estimated-wait: 2s
  max-queue-wait: 2s
  retry-after-seconds: 2
  throughput-window: 30s
  per-function-min-samples: 50
\end{lstlisting}

\paragraph{Che cosa governano i blocchi nuovi.}
\code{nanofaas.scheduler.strategy} sceglie la strategia su cui parte il motore
composto (\code{per-function} o \code{shared-queue}, \cref{sec:09-switch}); il
valore vuoto non \`e una terza strategia ma ``nessuna scelta esplicita'' e il
motore deriva la mappatura dai moduli di coda presenti nell'artefatto. Un
identificatore che nessuna strategia fornisce fa fallire l'avvio. Il cambio a
caldo passa dall'endpoint di amministrazione (\cref{cap:A-api}); l'ordine di
grandezza del budget di preparazione \`e una costante del motore (50\,ms), non
una chiave, e le soglie di verifica sono congelate in
\code{docs/experiments/scheduler-switching-2026-09/budgets.json}.
\code{nanofaas.retry.*} governa il backoff tra i tentativi
(\cref{sec:06-riprova}). \code{nanofaas.execution-store.*} governa la
ritenzione degli esiti: \code{max-outcome-bytes} \`e l'unico limite che la
cache degli esiti fa rispettare, mentre \code{max-outcomes} serve solo a
derivarne il default; \code{max-keys} limita le chiavi di idempotenza, e a
budget esaurito una \emph{nuova} ammissione con chiave \`e rifiutata con 429
prima del dispatch. \code{nanofaas.http-client.*} descrive il pool di
connessioni verso le funzioni, per destinazione concreta e non in aggregato.
\code{nanofaas.registry.path} \`e il file su cui persiste il catalogo delle
funzioni. \code{nanofaas.deployment.wakeup.*} regola l'attesa dell'avvio di una
funzione scalata a zero.

\subsection{Tetti di capacit\`a di invocazione}

Il blocco \code{nanofaas.invocation-capacity} fissa i tetti di ci\`o che pu\`o
essere vivo contemporaneamente: sono limiti di occupazione, non limiti di
frequenza, e non promettono un RSS. I default sono calibrati sul minimo
supportato dal chart, un container da 512\,MiB.

\begin{center}
\small
\begin{adjustbox}{max width=\textwidth}
\begin{tabular}{@{}lrr@{}}
\toprule
\textbf{Chiave} (\code{nanofaas.invocation-capacity.}) & \textbf{Globale} &
\textbf{Per funzione} \\
\midrule
\code{executions-*} & 4096 & 512 \\
\code{canonical-input-bytes-*} & 134\,217\,728 & 33\,554\,432 \\
\code{physical-input-copy-bytes-*} & 67\,108\,864 & 16\,777\,216 \\
\code{waiters-*} & 8192 & 1024 \\
\midrule
\code{ingress-body-bytes} & \multicolumn{2}{r}{1\,048\,576 (1\,MiB)} \\
\code{max-input-references} & \multicolumn{2}{r}{64} \\
\code{retained-input-max-depth} & \multicolumn{2}{r}{32} \\
\code{retained-input-max-container-entries} & \multicolumn{2}{r}{16\,384} \\
\code{retained-input-max-visited-nodes} & \multicolumn{2}{r}{65\,536} \\
\code{retained-input-max-bytes-per-execution} & \multicolumn{2}{r}{1\,048\,576} \\
\bottomrule
\end{tabular}
\end{adjustbox}
\end{center}

\section{Chiavi dei moduli non presenti nel file di default}

Ogni modulo dichiara in \code{module.properties} se \`e abilitato di default
(\code{defaultEnabled}). Lo \`e ogni modulo tranne i due provider di deployment
locali, \code{container-deployment-provider} e
\code{containerd-deployment-provider}: \code{sync-queue} \`e quindi attivo in
ogni build che lo include, senza bisogno di impostare una chiave. Il profilo di
ammissione (\code{nanofaas.admission.profile}: \emph{direct},
\emph{function-queue} o \emph{sync-queue}), se non \`e dichiarato, \`e derivato
dalle strategie presenti: \code{per-function} d\`a \emph{function-queue},
altrimenti \code{shared-queue} d\`a \emph{sync-queue}, altrimenti \emph{direct}.

\begin{center}
\small
\begin{adjustbox}{max width=\textwidth}
\begin{tabular}{@{}llp{5.4cm}@{}}
\toprule
\textbf{Chiave} & \textbf{Default} & \textbf{Modulo} \\
\midrule
\code{nanofaas.concurrency-control.poll-interval-ms} & 5000 &
\code{concurrency-control} \\
\code{nanofaas.concurrency-control.default-target-in-flight-per-pod} & 2 & \\
\code{nanofaas.concurrency-control.total-budget} & $\max(8, 4 \times$ core$)$ & \\
\midrule
\code{nanofaas.offload.enabled} & \code{true} & \code{offload} \\
\code{nanofaas.offload.target-url} & --- & \\
\code{nanofaas.offload.pressure-enabled} & \code{true} & \\
\midrule
\code{nanofaas.container-local.runtime-adapter} & \code{docker} &
\code{container-deployment-provider} \\
\code{nanofaas.container-local.bind-host} & \code{127.0.0.1} & \\
\code{nanofaas.container-local.readiness-timeout} & 20\,s & \\
\code{nanofaas.container-local.readiness-poll-interval} & 250\,ms & \\
\code{nanofaas.container-local.network-name} & --- & \\
\code{nanofaas.container-local.cpuset} & --- & vincola le funzioni agli stessi
core (\cref{cap:15-deployment-provider}) \\
\midrule
\code{nanofaas.containerd.socket-path} & \code{\$XDG\_RUNTIME\_DIR/containerd/containerd.sock} &
\code{containerd-deployment-provider} \\
\code{nanofaas.containerd.namespace} & \code{nanofaas} & \\
\code{nanofaas.containerd.runtime-binary} & \code{crun} & \\
\code{nanofaas.containerd.snapshotter} & \code{native} & \\
\code{nanofaas.containerd.network-name} & \code{nanofaas} & \\
\code{nanofaas.containerd.bind-host} & \code{127.0.0.1} & \\
\code{nanofaas.containerd.cgroups-path}, \code{cpuset}, \code{callback-url}, \code{state-directory}, \code{cni-*}, \code{readiness-*}, \code{*-timeout} & vedi codice & \\
\bottomrule
\end{tabular}
\end{adjustbox}
\end{center}

\section{Configurazione per funzione}

Tutti i campi di \code{FunctionSpec}. \code{null} significa ``non specificato'' e
attiva il default di piattaforma.

\begin{lstlisting}[language=yamlnf]
name: <obbligatorio, non vuoto>
image: <obbligatorio, non vuoto>
command: [...]                  # default []
env: { ... }                    # default {}
resources:
  requests: { cpu: <positivo>, memoryMiB: <positivo> }
  limits:   { cpu: <positivo>, memoryMiB: <positivo> }   # richiesta <= limite
timeoutMs: <>= 1>
concurrency: <>= 1>
queueSize: <>= 1>
maxRetries: <>= 0>
endpointUrl: <obbligatorio per EXTERNAL>
executionMode: LOCAL | EXTERNAL | DEPLOYMENT      # default DEPLOYMENT
runtimeMode: HTTP | STDIO | FILE
runtimeCommand: <comando del watchdog>
imagePullSecrets: [...]
scalingConfig:
  strategy: HPA | INTERNAL | NONE                 # default INTERNAL
  minReplicas: <>= 0>                             # default 1
  maxReplicas: <>= 1>                             # default 10
  metrics:
    - type: queue_depth | in_flight | rps
      target: "<numero>"
  concurrencyControl:
    mode: FIXED | STATIC_PER_POD | ADAPTIVE_PER_POD | BUDGETED | SOJOURN
    targetInFlightPerPod: 2
    minTargetInFlightPerPod: 1
    maxTargetInFlightPerPod: 8
    upscaleCooldownMs: 30000
    downscaleCooldownMs: 60000
    highLoadThreshold: 0.5
    lowLoadThreshold: 0.15
    targetLatencyMs: 250       # servizio per BUDGETED, end-to-end per SOJOURN
    weight: 1.0
offload:
  enabled: <true|false>
  targetUrl: <url>
  mode: pressure | always
\end{lstlisting}

\section{Costanti non configurabili}

Valori fissati nel codice, elencati perch\'e sono i candidati naturali per un
esperimento che voglia variarli.

\begin{center}
\small
\begin{adjustbox}{max width=\textwidth}
\begin{tabular}{@{}llp{5cm}@{}}
\toprule
\textbf{Costante} & \textbf{Valore} & \textbf{Dove} \\
\midrule
Batch massimo per funzione (strategia \code{per-function}) & 2 & \code{async-queue} \\
Finestra di scansione della coda condivisa (\code{shared-queue}) & 64 & \code{sync-queue} \\
Budget di preparazione di un cambio di scheduler & 50\,ms & motore (\code{SchedulerEngine}) \\
Cooldown di scale-up / scale-down & 30\,s / 60\,s & \code{autoscaler} \\
Finestra di inferenza del cold start & 30\,s & \code{autoscaler} \\
Deriva della baseline & $1{,}002$ per tick & \code{concurrency-control} \\
Decadimento degli estremi & $1{,}01$ per tick & \code{concurrency-control} \\
Pavimento del gradiente & $0{,}5$ & \code{concurrency-control} \\
Soglia di guardia della coda & $0{,}7$ & \code{concurrency-control} \\
Orizzonte di drenaggio & 5\,s (min 1\,s) & \code{concurrency-control} \\
Holt: $\alpha$, $\beta$, $\phi$ & $0{,}5$, $0{,}3$, $0{,}8$ &
\code{concurrency-control} \\
Margine di timeout dell'offload & 50\,ms & \code{offload} \\
\bottomrule
\end{tabular}
\end{adjustbox}
\end{center}

\section{Selettori di build}

\begin{center}
\small
\begin{adjustbox}{max width=\textwidth}
\begin{tabular}{@{}lp{7.6cm}@{}}
\toprule
\textbf{Selettore} & \textbf{Effetto} \\
\midrule
\code{-PcontrolPlaneModules} & \code{none}, \code{all} (default) o lista; anche
\code{NANOFAAS\_CONTROL\_PLANE\_MODULES} \\
\code{-PnativeOptimization} & \code{3} (default) o \code{s} \\
\code{-PnativeGc} & \code{G1} richiede Oracle GraalVM; anche
\code{NATIVE\_GC} \\
\code{-PnativeMonitoring} & per esempio \code{jfr,jvmstat} \\
\code{-PnativeBuildMemory} & heap del \emph{compilatore}, non dell'immagine \\
\code{-PnativeParallelism} & lavoratori del compilatore \\
\bottomrule
\end{tabular}
\end{adjustbox}
\end{center}
